% ===================================================================
%  QUADRO N.º 1 — A Escola. — O Liceu. — A Sala de aula.
%
%  Ceci n'est PAS une transcription : c'est une traduction, faite en
%  2026, en portugais d'aujourd'hui. D'ou trois differences avec
%  texto/io et texto/fr, et trois seulement :
%
%   * pas de \nl ni de \cc — la traduction ne suit aucune ligne
%     imprimee, et les lignes de ce fichier ne sont que du confort ;
%   * pas de \begin{VUpage} — elle n'a ni feuillet ni folio, et la
%     page de lecture ne lui en affiche pas ;
%   * le gras se pose sur le TERME seul, ponctuation dehors.
%
%  Tout le reste est identique, et doit l'etre : les cles « %%K » sont
%  celles de texto/io, macro pour macro pour l'apparat, et les renvois
%  \textsuperscript{(n)} visent les MEMES objets de la planche — ce
%  sont eux qui ouvrent les gros plans.
%
%  L'ido est le texte de base ; le francais sert de controle, et
%  l'anglais et l'espagnol de calibrage : les traductions doivent
%  s'accorder entre elles autant qu'avec la source.
%
%  LE GRAS NE SE POSE PAS SUR UN MOT-OUTIL. Le fac-simile ido le fait
%  par accident — « \VUgras{Ica} docisto » — et la page de lecture l'en
%  corrige ; la traduction, elle, n'a pas a repeter l'accident.
%
%  LES NOMS DE PERSONNE SONT NATURALISES, comme dans les sept autres
%  colonnes : Ioannes est Joao, Iakobus est Jaime. Seuls le chien
%  Medor et l'ane Aliboron gardent leur nom.
%
%  L'ORTHOGRAPHE EST CELLE DE L'ACCORD DE 1990, et le vocabulaire
%  celui du Bresil la ou les deux rives divergent — c'est la ou le
%  portugais compte le plus de locuteurs. « Antonio » prend donc son
%  accent circonflexe, et l'on ecrit « trem » et non « comboio ».
%
%  LA FAUTE N'ETAIT PAS DANS LE PORTUGAIS, ELLE ETAIT DANS LA MANIERE
%  DE TRADUIRE. Vingt-deux alineas de cette colonne ouvraient par
%  « eis », decalque du « voici » que le fac-simile francais place en
%  tete de presque chaque objet nouveau. « Eis » n'est pas fautif : il
%  est simplement du registre ecrit soutenu, et un livret de 2026 qui
%  montre une planche murale ne l'emploie pas vingt-deux fois. On l'a
%  remplace cas par cas — « ao seu lado esta », « numa taca grande ha »
%  — et non par une formule unique.
%
%  Les MEMES blocs portaient « he aqui » en espagnol, « ty » en
%  suedois, « for » en anglais : la tournure vieillie voyage avec le
%  bloc, d'ou qu'on le traduise.
% ===================================================================

%%K t01-apar-0 apar
\VUsaut{2.0mm}
\VUcentre{12.6pt}{120}{FOLHETO EXPLICATIVO}
\VUsaut{3.2mm}
\VUcentre{8.4pt}{160}{DOS}
\VUsaut{2.6mm}
\VUtitre{15.6pt}{0}{Quadros Auxiliares Delmas}
\VUsaut{3.4mm}
\VUornamento{9pt}
\VUsaut{5.0mm}
\VUcentre{13.2pt}{90}{PRIMEIRA SÉRIE}
\VUsaut{3.0mm}
\VUfilet{16mm}
\VUsaut{5.6mm}

%%K t01-apar-1 apar
\VUtitre{15.0pt}{60}{QUADRO N.º 1}
\VUsaut{1.4mm}
\VUtitre{11.4pt}{0}{A Escola. --- O Liceu. --- A Sala de aula.}
\VUsaut{2.2mm}
\VUfilet{20mm}
\VUsaut{6.4mm}

%%K t01-tit-1 sub
\VUpk{10.2pt}{40}{Descrição geral.}
\VUsaut{3.0mm}

%%K t01-01-1 p
1. --- Este quadro representa uma \VUgras{sala de aula} de \VUgras{liceu}.
Nesta sala há oito \VUgras{mesas} \textsuperscript{(18)} e oito \VUgras{bancos}
\textsuperscript{(32)}. Em cada banco sentam-se três alunos. O
\VUgras{professor} \textsuperscript{(7)} está sentado numa \VUgras{cadeira}
\textsuperscript{(8)} na sua \VUgras{cátedra} \textsuperscript{(6)}.

%%K t01-02-1 p
2. --- À esquerda da cátedra há um \VUgras{armário} com portas de vidro
\textsuperscript{(15)}. À direita está o \VUgras{quadro-negro}
\textsuperscript{(1)} com o \VUgras{apagador} \textsuperscript{(2)} e o \VUgras{giz}
\textsuperscript{(3)}.

%%K t01-02-2 p
Por cima da cátedra estão penduradas umas \VUgras{estampas murais}
\textsuperscript{(9, 11, 12)} e um \VUgras{mapa} \textsuperscript{(10)}.

%%K t01-03-1 p
3. --- Nem todos os alunos estão no seu \VUgras{lugar} \textsuperscript{(17)}. João
\textsuperscript{(4)} está de pé junto ao quadro-negro; segura o apagador e
apaga as \VUgras{figuras de geometria}.

%%K t01-03-2 p
Pedro está perto da cátedra; recolhe os \VUgras{exercícios escritos}
\textsuperscript{(5)} e leva-os ao professor.

%%K t01-04-1 p
4. --- \VUgras{Este} professor é o professor de \VUgras{línguas modernas}.
Ensina os alunos a falar, ler e escrever línguas estrangeiras
\VUgras{(alemão, inglês, espanhol, italiano, russo} ou \VUgras{francês)}.

%%K t01-05-1 p
5. --- A sala de aula é muito ampla e muito clara. Ao fundo há uma
\VUgras{janela} larga com duas \VUgras{bandeiras} \textsuperscript{(78)} e uma
\VUgras{porta envidraçada} para arejá-la e iluminá-la.

%%K t01-05-2 p
Esta sala não é fria. No meio, para aquecê-la, há um \VUgras{fogão}
\textsuperscript{(46)} muito bom com o seu \VUgras{tubo} \textsuperscript{(45)}.

%%K t01-05-3 p
No tabique estão fixados uns \VUgras{cabides} \textsuperscript{(58)}; deles pendem
os gorros e os casacos dos alunos.

%%K t01-tit-2 sub
\VUsaut{5.2mm}
\VUpk{10.2pt}{40}{O pátio de recreio.}
\VUsaut{3.6mm}

%%K t01-06-1 p
6. --- O \VUgras{relógio} \textsuperscript{(63)} marca nove e cinco.

%%K t01-06-2 p
O \VUgras{recreio} \textsuperscript{(76)} acaba de terminar e a aula recomeça. O
recreio faz-se no \VUgras{pátio} \textsuperscript{(75)}. Vemos alguns alunos que
brincam. Um \VUgras{monitor} \textsuperscript{(74)} vigia-os.

%%K t01-07-1 p
7. --- No pátio vemos ainda outras pessoas: o \VUgras{inspetor}
\textsuperscript{(73)}, o \VUgras{diretor} \textsuperscript{(71)} e o
\VUgras{chefe de estudos} \textsuperscript{(72)}.

%%K t01-07-2 p
O inspetor inspeciona as escolas, o diretor dirige o liceu, o chefe de
estudos vigia os estudos.

%%K t01-07-3 p
A quarta personagem é o \VUgras{porteiro} \textsuperscript{(77)}. Leva debaixo do
braço um registo para anotar os ausentes.

%%K t01-08-1 p
8. --- Ao fundo do pátio estão os \VUgras{aparelhos de ginástica}
\textsuperscript{(70)} e a oficina de \VUgras{trabalhos manuais}
\textsuperscript{(69)}.

%%K t01-08-2 p
À direita do quadro entrevemos as \VUgras{salas} de \VUgras{música} e de
\VUgras{desenho}.

%%K t01-noto-1 noto
\VUnotes{143.2mm}{%
(*) Para os \textit{nomes próprios} aconselha-se transcrevê-los também
na sua forma latina. É o único modo de evitar variantes sem conta. De
resto, como escolher entre \textit{Giovanni, Jean, Johann, Jan, Hans,}
\textit{John, Ivan,} etc.? Será preciso criar um dicionário especial de
nomes próprios? Tomemos do latim o seu nominativo e digamos:
\VUgras{Ioannes, Ioanna; Iulius, Iulia; Iakobus; Andreas; Lukas.}
(Gramática completa).%
}

%%K t01-tit-3 sub
\VUpk{10.2pt}{40}{Descrição pormenorizada.}
\VUsaut{2.4mm}

%%K t01-tit-4 sub
\VUpk{10.2pt}{40}{Os bons alunos.}
\VUsaut{3.0mm}

%%K t01-09-1 p
9. --- Os alunos do primeiro banco à direita da cátedra são
\VUgras{Henrique} (*), \VUgras{Estêvão} e \VUgras{Carlos}.

%%K t01-09-2 p
Henrique é muito atento e trabalhador; escuta o professor. Sobre a sua
mesa tem um \VUgras{caderno} \textsuperscript{(28)}, um \VUgras{livro}
\textsuperscript{(29)}, um \VUgras{dicionário} \textsuperscript{(30)} e um
\VUgras{estojo} \textsuperscript{(31)}. O seu livro tem muitas \VUgras{páginas};
cada capítulo contém vários \VUgras{parágrafos} e cada \VUgras{linha}
muitas \VUgras{palavras}.

%%K t01-09-4 p
O seu vizinho Estêvão \textsuperscript{(24)} é aplicado. Anota no seu caderno a
lição para a próxima vez. Tem boa \VUgras{letra}. O seu livro está
fechado. Só vemos a \VUgras{capa} \textsuperscript{(27)}. Ao seu lado está o seu
\VUgras{compasso} \textsuperscript{(26)}.

%%K t01-09-5 p
Carlos \textsuperscript{(23)} segura o seu livro na mão; abre-o para rever a
lição. Como todos os alunos, tem diante de si um \VUgras{tinteiro}
\textsuperscript{(25)}. É obediente.

%%K t01-10-1 p
10. --- Na mesa ao lado estão \VUgras{Luís, Antônio} e \VUgras{Paulo}. Luís
recita a sua \VUgras{lição} \textsuperscript{(22)} e responde às \VUgras{perguntas}
do professor. Antônio \textsuperscript{(21)} está muito distraído; às vezes o
professor até o castiga. Paulo \textsuperscript{(20)} é bom companheiro, amável
e prestativo. É muito atento.

%%K t01-tit-5 sub
\VUsaut{4.2mm}
\VUpk{10.2pt}{40}{Os maus alunos.}
\VUsaut{3.2mm}

%%K t01-11-1 p
11. --- Atrás de Henrique está \VUgras{Jorge} \textsuperscript{(38)}. Não é mau
rapaz, mas é \VUgras{falador}, distraído e irrequieto. A sua
\VUgras{leviandade} e a sua \VUgras{desobediência} causam \VUgras{desordem}
durante a aula, e muitas vezes merece uma severa \VUgras{advertência}. O
seu \VUgras{caderno de rascunho} \textsuperscript{(35)} está sujo, enche-o de
\VUgras{borrões de tinta} com a sua \VUgras{caneta} \textsuperscript{(34)}, e por
isso usa sempre a sua \VUgras{borracha} \textsuperscript{(36)}; a sua
\VUgras{pasta} \textsuperscript{(33)} está rota, porque é muito descuidado. Por
que traz o seu \VUgras{porta-lápis} \textsuperscript{(37)} à aula de línguas
modernas?

%%K t01-11-2 p
O seu vizinho Edmundo \textsuperscript{(39)} não gosta dele. É verdade que é um
tanto preguiçoso, mas é calado e fica quieto. Não conversa; só que, em
vez de atender, prefere ler os \VUgras{contos} do seu \VUgras{livro de
leitura}.

%%K t01-11-3 p
O terceiro aluno deste banco é \VUgras{Ludovico} \textsuperscript{(40)}. É
aplicado. Escreve numa \VUgras{folha de papel} branca. Diante de si
deixou o seu \VUgras{mata-borrão} \textsuperscript{(41)}.

%%K t01-12-1 p
12. --- Os alunos do segundo banco, à esquerda do professor, são muito
jovens. O primeiro, o pequeno \VUgras{Emílio}, ainda não sabe escrever
bem; traz a sua \VUgras{lousa} \textsuperscript{(42)}, o seu \VUgras{lápis de
lousa} e a sua \VUgras{esponja} \textsuperscript{(43)}.

%%K t01-12-2 p
Ao seu lado estão \VUgras{Augusto} e \VUgras{Bernardo} \textsuperscript{(44)}. Este
último trabalha bem, mas o seu \VUgras{aspeto} é muito descuidado.

%%K t01-13-1 p
13. --- Atrás deles há três \VUgras{preguiçosos}. \VUgras{Alberto} não aprende
as lições. \VUgras{Jaime} \textsuperscript{(47)} dorme em vez de escutar. A sua
\VUgras{preguiça} acarreta-lhe \VUgras{repreensões} e até \VUgras{castigos}.
Por fim, \VUgras{Cristiano} \textsuperscript{(48)} é frívolo demais.
\VUgras{Comporta-se} mal e não faz nenhum esforço por emendar-se.

%%K t01-14-1 p
14. --- Os seus vizinhos, pelo contrário, comportam-se bem.
\VUgras{Ernesto} \textsuperscript{(51)} ama a ordem e a regularidade; usa sempre
a sua \VUgras{régua} \textsuperscript{(50)} e o seu \VUgras{lápis}
\textsuperscript{(49)}. É \VUgras{externo}.

%%K t01-14-2 p
\VUgras{Francisco} \textsuperscript{(52)} é \VUgras{interno}; usa o \VUgras{uniforme}
do liceu; ali come e ali dorme. É um bom \VUgras{companheiro}.

%%K t01-14-3 p
Maurício apara o seu \VUgras{lápis} com o seu \VUgras{canivete}
\textsuperscript{(56)}; é muito cuidadoso.

%%K t01-14-4 p
Traz todos os seus livros, a sua \VUgras{gramática} \textsuperscript{(55)}, o seu
\VUgras{caderno de apontamentos} \textsuperscript{(54)} e até uma \VUgras{pasta de
escrivaninha} \textsuperscript{(53)}. A sua \VUgras{mochila escolar}
\textsuperscript{(57)} está no chão atrás dele.

%%K t01-14-5 p
As duas mesas do fundo da sala são ocupadas por \VUgras{bons alunos}
\textsuperscript{(19)}.

%%K t01-tit-6 sub
\VUsaut{4.6mm}
\VUpk{10.2pt}{40}{Desenho e música.}
\VUsaut{3.4mm}

%%K t01-15-1 p
15. --- Na \VUgras{sala de desenho} há bonitos \VUgras{modelos} pendurados na
parede e um \VUgras{candeeiro} \textsuperscript{(62)} com \VUgras{quebra-luz},
suspenso do teto. O \VUgras{professor} \textsuperscript{(60)} é muito paciente;
corrige os \VUgras{desenhos} \textsuperscript{(61)} dos alunos e ensina-lhes a
desenhar um \VUgras{busto} \textsuperscript{(59)} ao natural.

%%K t01-16-1 p
16. --- A \VUgras{sala de música} parece muito interessante. Ali aprende-se
a ler a \VUgras{música}, a solfejar as \VUgras{notas da escala}
\textsuperscript{(67)} escritas na \VUgras{pauta} (dó, ré, mi, fá, sol, lá,
si\,=\,C. D. E. F. G. A. B.), a reconhecer as \VUgras{claves} e a cantar
belas \VUgras{melodias}. As partituras põem-se na \VUgras{estante}
\textsuperscript{(65)}. Um dos alunos \VUgras{toca piano} \textsuperscript{(64)} e o
\VUgras{professor de música} \textsuperscript{(66)} \VUgras{marca o compasso}
\textsuperscript{(68)}.

%%K t01-17-1 p
17. --- Entrevemos um canto da \VUgras{oficina} de \VUgras{trabalhos manuais}
\textsuperscript{(69)}, onde as crianças podem aprender a aplainar a madeira e
a limar o ferro. Não a há em todos os liceus.

%%K t01-tit-7 sub
\VUsaut{5.0mm}
\VUpk{10.2pt}{40}{A sala de ciências.}
\VUsaut{3.6mm}

%%K t01-18-1 p
18. --- O professor de matemática ensina aos alunos \VUgras{geometria,}
\VUgras{aritmética, cálculo} e o \VUgras{sistema métrico}. Vemos no
quadro-negro as principais figuras de geometria: um \VUgras{círculo}
\textsuperscript{(a)}, um \VUgras{quadrado} \textsuperscript{(b)}, um \VUgras{triângulo}
\textsuperscript{(c)}, uma \VUgras{linha} \textsuperscript{(d)}.

%%K t01-18-2 p
Vemos também uma \VUgras{operação}; não é uma \VUgras{soma}, nem uma
\VUgras{subtração}, nem uma \VUgras{multiplicação}: é uma \VUgras{divisão}
\textsuperscript{(e)}.

%%K t01-19-1 p
19. --- Na estampa do sistema métrico distinguimos: um \VUgras{metro}
\textsuperscript{(a)}, uma \VUgras{balança} \textsuperscript{(b)} e os seus
\VUgras{pesos}, um \VUgras{litro} \textsuperscript{(e)}, um \VUgras{decalitro}
\textsuperscript{(f)}, um \VUgras{decilitro} e um \VUgras{metro cúbico}
\textsuperscript{(d)}.

%%K t01-19-2 p
Os alunos estudam também \VUgras{ciências naturais} \textsuperscript{(11)} ---
zoologia \textsuperscript{(a)}, botânica \textsuperscript{(b)}, geologia
\textsuperscript{(c)} --- e \VUgras{história} \textsuperscript{(12)}.

%%K t01-19-3 p
No armário estão os instrumentos de \VUgras{física} \textsuperscript{(15)} e de
\VUgras{química} \textsuperscript{(16)}.

%%K t01-19-4 p
Por cima do armário há: uma \VUgras{esfera} \textsuperscript{(14)}, um \VUgras{cone}
e um \VUgras{globo terrestre} \textsuperscript{(13)}.

%%K t01-19-5 p
Por cima das estampas pende um \VUgras{mapa} que representa as cinco
partes do mundo: \VUgras{Europa} \textsuperscript{(g)}, \VUgras{Ásia}
\textsuperscript{(e)}, \VUgras{África} \textsuperscript{(f)}, \VUgras{América}
\textsuperscript{(ab)} e \VUgras{Oceania} \textsuperscript{(h)}, e os dois grandes
oceanos, o \VUgras{Pacífico} \textsuperscript{(c)} e o \VUgras{Atlântico}
\textsuperscript{(d)}.
\VUsaut{6.0mm}
\VUfilet{22mm}
